\section{文字の扱いの変遷}
\label{sec_history}
表\ref{tab_history}に，定義書の各版における文字の扱いをまとめる．
文字の扱いは，読める文字の有無を表す文字ありと，題として扱う文字を表すお題文の2つの項目に分かれて改訂されてきた．

\begin{table*}[t]
    \centering
    \caption{定義書の各版における文字の扱い}
    \label{tab_history}
    \small
    \begin{tabular}{>{\centering\arraybackslash}p{0.07\textwidth}>{\raggedright\arraybackslash}p{0.36\textwidth}>{\raggedright\arraybackslash}p{0.25\textwidth}>{\raggedright\arraybackslash}p{0.23\textwidth}}
        \toprule
        版 & お題文の扱い & 文字ありの扱い & 改訂の根拠 \\
        \midrule
        第1版 & 問い，指示，題目の句，状況の設定，名前の列挙．字幕や看板など元からある文字は除く & お題文ありと，その他の文字ありを別の属性とする & -- \\
        第2版 & テロップ，字幕，看板など投稿者が重ねたのではない文は，問いの形でもお題文としない & 同上 & 教師のお題文ありの F1 が約0.63 \\
        第3版 & 形式のお題文のみで表す & お題文ありとその他の文字ありを文字ありに統合．ロゴや透かしは除く & 2項目の境界の誤りが多く，下流の処理は文字を読むかのみを要する \\
        第4版 & テンプレートや無地の背景に書いた主題の文を含めると明記 & ロゴ，透かし，衣服の文字など，読める文字全般とする & 誤りの多くがロゴや小さな文字の扱い \\
        第5版 & お題文ありを空欄やカルタ記号と独立の属性とする．空欄を含む文そのもの，元の素材の文字，吹き出しの台詞を除く & 第4版と同じ & 穴埋めやカルタの画像にも問いの文がある \\
        第6版（提案） & 文字の層と文の形の属性から規則で導く．誰が書いたかを判定に用いない & 層のどれかがあれば文字ありとする & 意図に基づく判断の揺れ \\
        \bottomrule
    \end{tabular}
\end{table*}

文字ありは，改訂のたびに観測できる基準へ近づいた．
第1版と第2版では，お題文ありとその他の文字ありを別の属性としていたが，2つの項目の誤りの多くは，テロップを投稿者が重ねた文とみなすかの判断にあった．
下流の処理で必要なのは文字を読み取るかの判断であり，題であるかは形式のお題文で表せるため，第3版で2つの項目を文字ありに統合した．
第3版の文字ありはロゴや透かしを除いていたが，教師はこれらを含めて回答し，定義と教師の判定がずれていた．
そこで第4版では，読める文字があれば文字ありとした．
この改訂では，ロゴや衣服の文字を問う別の短い質問を教師に加え，カルタ記号のない画像に限ってその回答を文字ありに加えた．
その結果，実データの目視ラベルに対する教師の文字ありの F1 は0.954から0.971に上がった．
第4版の目視ラベルに対して，パッケージの文字ありは実データで適合率0.94，再現率0.99であり，文字ありの判定に大きな課題は残っていない．

一方，お題文は，投稿者が回答者に向けて書いた文という意図に基づく定義を第1版から保ち，境界の規則を版ごとに書き足してきた．
第2版では，元からある文字の範囲をテロップや字幕にまで広げた．
この改訂により，教師のお題文ありの F1 は約0.63から0.79に上がった．
第5版では，空欄やカルタ記号のある画像にも投稿者が書いた問いがあることから，お題文ありを形式から独立した属性とした．
このとき，話者を付けた台本，分類の名札だけの見出し，書き写した既存の歌詞，テンプレートの定型の見出しをお題文から除く規則を加えた．
これらの規則は，いずれも文字の出所や書き手の意図を判断の根拠としている．
